%% kapa: exam=2 round=30 year=2019 period=3 subject=compensation_law subject_ko=감정평가 및 보상법규 seq=1 number=1 points=40
%% source: Q-Net 감정평가사 기출문제 (cst00309) artlSeq=5205595 file=raw/1008/5205595/감정평가 및 보상법규.pdf
%% conversion: kordoc 4.24.51 -> tools/gen2.py (auto); review=proofread-gpt-terra; flags=none
\begin{longquestion}{【문제 1】}
관할 A시장은「부동산 가격공시에 관한 법률」에 따라 甲 소유의 토지에 대해 공시기준일을 2018. 1. 1.로 한 개별공시지가를 2018. 6. 28. 결정ㆍ공시하고(‘당초 공시지가’) 甲에게 개별통지하였으나, 이는 토지가격비준표의 적용에 오류가 있는 것이었다. 이후 甲 소유의 토지를 포함한 지역 일대에 개발사업이 시행되면서 관련법에 의한 부담금 부과의 대상이 된 甲의 토지에 대해 A시장은 2018. 8. 3. 당초 공시지가에 근거하여 甲에게 부담금을 부과하였다. 한편 甲 소유 토지에 대한 당초 공시지가에 이의가 있는 인근 주민 乙은 이의신청기간이 도과한 2018. 8. 10. A시장에게 이의를 신청하였고, A시장은 甲 소유 토지에 대한 당초 공시지가를 결정할 때 토지가격비준표의 적용에 오류가 있었음을 이유로「부동산 가격공시에 관한 법률」제12조 및 같은법 시행령 제23조제1항에 따라 개별공시지가를 감액하는 정정을 하였고, 정정된 공시지가는 2018. 9. 7. 甲에게 통지되었다. 다음 물음에 답하시오. (아래 설문은 각각 별개의 독립된 상황임) (40점) \par
\medskip\noindent 물음 1) 甲은 정정된 공시지가에 대해 2018. 10. 22. 취소소송을 제기하였다. 甲의 소송은 적법한가? (15점) \par
\medskip\noindent 물음 2) 甲은 이의신청기간이 도과한 후에 이루어진 A시장의 개별공시지가 정정처분은 위법하다고 주장한다. 甲의 주장은 타당한가? (10점) \par
\medskip\noindent 물음 3) 만약, A시장이 당초 공시지가에 근거하여 甲에게 부담금을 부과한 것이 위법한 것이더라도, 이후 A시장이 토지가격비준표를 제대로 적용하여 정정한 개별공시지가가 당초 공시지가와 동일하게 산정되었다면, 甲에 대한 부담금 부과의 하자는 치유되는가? (15점) \par
\begin{framedtext}
참조 조문 \par
<부동산 가격공시에 관한 법률> \par
제11조(개별공시지가에 대한 이의신청) ① 개별공시지가에 이의가 있는 자는 그 결정ㆍ공시일부터 30일 이내에 서면으로 시장ㆍ군수 또는 구청장에게 이의를 신청할 수 있다. \par
②～③ 생략 \par
제12조(개별공시지가의 정정) 시장ㆍ군수 또는 구청장은 개별공시지가에 틀린 계산, 오기, 표준지 선정의 \par
착오, 그 밖에 대통령령으로 정하는 명백한 오류가 있음을 발견한 때에는 지체 없이 이를 정정하여야 한다. \par
<부동산 가격공시에 관한 법률 시행령> \par
제23조(개별공시지가의 정정사유) ① 법 제12조에서 "대통령령으로 정하는 명백한 오류"란 다음 각 호의 \par
어느 하나에 해당하는 경우를 말한다. \par
1. 법 제10조에 따른 공시절차를 완전하게 이행하지 아니한 경우 \par
2. 용도지역ㆍ용도지구 등 토지가격에 영향을 미치는 주요 요인의 조사를 잘못한 경우 \par
3. 토지가격비준표의 적용에 오류가 있는 경우 \par
② 생략 \par
\end{framedtext}
\end{longquestion}
